\section*{Concluding Remarks}

Across a range of PDEs, this work identifies a basic limitation of purely empirical neural operator training: ``learning by experience'' incurs rapidly increasing data demands as the number of independent input DOF grows. Sensitivity-constrained training ``enforces the relationships'' and introduces an instructional paradigm in which the model is informed not only of the outcomes but also of how inputs must shape those outcomes, analogous to ``enforcing known principles'' during training. This sensitivity-driven scaling fundamentally alters the learning economics of AI-infused PDE solvers, enabling robust forward and inverse inference for gridded, spatially distributed inputs that were previously impractical to handle. We provide a \textit{Theoretical Justification for Sensitivity-Constrained Training} (Appendix~\ref{sec_sup_method}) and mechanistic evidence for sensitivity-driven scaling. More broadly, these results highlight the value of incorporating structured relational information, such as adjoint information, into data-driven AI, for example in adjoint matching~\cite{pmlr-v267-havens25a}.


Because neural operators are differentiable and fast, the SC-NO pathway also enables highly refined PDE solvers to be used in hybrid AI--physics learning and optimization frameworks \cite{shen2023differentiable, StroferXiao2021, zong2026mathematics}. In such frameworks, neural networks are coupled to and trained together with physical models to learn missing relationships, correct structural bias, identify resolution-dependent parameterization, or discover unclear knowledge. Thousands of model runs are required for such learning, and only an efficient surrogate with high-fidelity sensitivity can support these critical missions. Finally, Jacobian calculation has long been viewed as prohibitively expensive \cite{cao2025derivative}. However, we argue that this cost must be evaluated in the context of the information benefits—the SC-NO achieves noticeably better accuracy at the same cost for most setups. This advantage arises, at its core, from the efficiency of reverse-mode differentiation (backpropagation), which computes gradients with respect to high-dimensional inputs at a cost comparable to a small constant multiple of a single forward evaluation. Reverse-mode differentiation has been extensively optimized in modern AI frameworks. In this work, all computations are performed on a small number of GPUs, indicating substantial room to scale to much larger systems, together with dimension reduction methods.


For tsunamis, the time between offshore detection and shoreline impact is typically only 20--120 minutes (depending on epicenter location), and each computational minute saved has real-world implications. With only $\sim$4.5 minutes (further reducible) of inversion and forward time demonstrated here, near--real-time, high-fidelity inference is finally possible. While neural operators are known for their extraordinary computational efficiency, we show that their previous scaling law is what has limited them so far. Sensitivity constraints and accurate dynamical inversion thus expand the frontier of computing for numerous time-critical applications, such as volcanic eruption forecasting \cite{rutarindwa2019dynamic}, flood inundation, chemical plume releases \cite{di2020kalman}, mine safety, and coastal storm surges.